\section{Overview of the Search Mechanism}
RAG (Retrieval-Augmented Generation) in this system combines two complementary search methods:

\begin{itemize}
    \item \textbf{BM25 Keyword Search}: Precise matching on category names and colors.
    \item \textbf{Marqo-FashionSigLIP Vector Search}: Semantic understanding, synonyms, and feature descriptions.
\end{itemize}

\section{Marqo-FashionSigLIP Embedding --- Semantic Token}

\subsection{Why FashionSigLIP Instead of the Original FashionCLIP?}
\textbf{Marqo-FashionSigLIP} (\texttt{Marqo/marqo-fashionSigLIP}) is a SigLIP-based model fine-tuned on fashion data, producing \textbf{768-dimensional vectors} (vs.\ 512-d from the original FashionCLIP). The model uses the \texttt{open\_clip} library instead of HuggingFace's CLIPModel, with better support for Apple MPS.

\begin{lstlisting}[caption={FashionEmbedder class --- indexing/build\_index.py}]
class FashionEmbedder:
    """Wraps Marqo-FashionSigLIP for image & text encoding."""
    MODEL_NAME = "hf-hub:Marqo/marqo-fashionSigLIP"

    def __init__(self, cache_dir=None):
        import open_clip
        self.model, _, self.preprocess = (
            open_clip.create_model_and_transforms(
                self.MODEL_NAME,
                cache_dir=str(cache_dir) if cache_dir else None,
        ))
        self.tokenizer = open_clip.get_tokenizer(self.MODEL_NAME)
        # Device selection: MPS (Apple Silicon) > CUDA > CPU
        if torch.backends.mps.is_available():
            self._device = "mps"
        elif torch.cuda.is_available():
            self._device = "cuda"
        else:
            self._device = "cpu"
        self.model = self.model.to(self._device).eval()

    def encode_image(self, image_path: str) -> list[float]:
        """Convert image to 768-d vector."""
        img = Image.open(image_path).convert("RGB")
        tensor = self.preprocess(img).unsqueeze(0).to(self._device)
        with torch.no_grad():
            feat = self.model.encode_image(tensor)
            feat = feat / feat.norm(dim=-1, keepdim=True)
        return feat[0].cpu().tolist()

    def encode_text(self, text: str) -> list[float]:
        """Convert text to 768-d vector."""
        tokens = self.tokenizer([text]).to(self._device)
        with torch.no_grad():
            feat = self.model.encode_text(tokens)
            feat = feat / feat.norm(dim=-1, keepdim=True)
        return feat[0].cpu().tolist()
\end{lstlisting}

\textbf{Task breakdown:}
\begin{itemize}
    \item \texttt{open\_clip}: A flexible embedding library supporting many CLIP/SigLIP variants.
    \item \textbf{L2 normalization}: $\hat{v} = v / \|v\|_2$ ensures dot product $\equiv$ cosine similarity.
    \item \textbf{MPS priority}: Leverages Apple Silicon GPU for faster encoding.
    \item \textbf{768-d}: A larger vector dimension than FashionCLIP (512-d), enabling richer semantic representation.
\end{itemize}

\section{BM25 --- Label and Color Search}
BM25 (Best Match 25) is applied on short content \texttt{``label color''} --- e.g., ``T-Shirt Navy Blue'' --- enabling precise keyword matching on category names and specific colors.

\begin{tcolorbox}[title=Method: Hybrid Search --- BM25 + FashionSigLIP]
\begin{itemize}
    \item BM25 excels at: precise category name and color matching (e.g., ``navy blue shirt'').
    \item FashionSigLIP Vector excels at: semantic understanding, synonyms, and feature descriptions (e.g., ``slim fit cotton casual'').
    \item RRF Fusion: combines results from both with $k = 60$, $\text{bm25\_weight}=1.0$, $\text{vec\_weight}=2.5$.
\end{itemize}
\end{tcolorbox}

\begin{lstlisting}[caption={BM25 content composition --- indexing/build\_index.py}]
def compose_bm25_content(label: str, color: str) -> str:
    """Create content for the BM25 index.
    Only uses label + color (NOT caption)
    for more precise keyword matching."""
    parts = []
    if label:
        parts.append(label.strip())
    if color:
        parts.append(color.strip())
    return " ".join(parts)
\end{lstlisting}

\section{RRF Fusion and Soft Relevance Scoring}

The Reciprocal Rank Fusion score is computed as in Equation~(\ref{eq:rrf}):
\begin{equation}
\label{eq:rrf}
\text{score}(d) \;=\; \sum_{r \in R} \frac{w_r}{k + \text{rank}_r(d) + 1}
\end{equation}

\begin{lstlisting}[caption={Reciprocal Rank Fusion --- search/fusion.py}]
@dataclass
class NodeWithScore:
    """Search result with full metadata."""
    image_id: str
    label: str = ""
    color: str = ""
    caption: str = ""
    image_path: str = ""
    bm25_content: str = ""
    score: float = 0.0

def reciprocal_rank_fusion(
    bm25_results: list[NodeWithScore],
    vec_results: list[NodeWithScore],
    k: int = 60,
    bm25_weight: float = 1.0,
    vec_weight: float = 2.5,
) -> list[NodeWithScore]:
    """Combine BM25 and Vector Search results via RRF."""
    fused: dict[str, dict] = {}
    for rank, node in enumerate(bm25_results):
        nid = node.image_id
        if nid not in fused:
            fused[nid] = {"node": node, "score": 0.0}
        fused[nid]["score"] += bm25_weight / (k + rank + 1)

    for rank, node in enumerate(vec_results):
        nid = node.image_id
        if nid not in fused:
            fused[nid] = {"node": node, "score": 0.0}
        fused[nid]["score"] += vec_weight / (k + rank + 1)

    sorted_items = sorted(
        fused.values(), key=lambda x: x["score"], reverse=True)
    for item in sorted_items:
        item["node"].score = item["score"]
    return [item["node"] for item in sorted_items]
\end{lstlisting}

\textbf{Task breakdown:}
\begin{itemize}
    \item \textbf{RRF}: Avoids dependence on absolute scores from individual retrievers (which use different scales). Each result is ranked by position using the formula in Eq.~(\ref{eq:rrf}).
    \item \textbf{vec\_weight = 2.5}: Vector search receives a higher weight because semantic matching is richer than keyword matching.
    \item \textbf{Soft filter}: \texttt{soft\_relevance\_filter()} uses RapidFuzz for fuzzy color/category comparison (threshold = 60), without completely excluding results that do not match exactly.
\end{itemize}

% ===========================================================================
% CHAPTER 4: POSTGRESQL INTEGRATION
% ===========================================================================
